% ECONOMICS_PROBLEM: {"id": "F33-118", "title": "Global financial safety net", "jel_code": "F33", "source_id": "OP-0339"}
% !TeX program = xelatex
\documentclass[11pt,a4paper]{article}
\usepackage[margin=23mm]{geometry}
\usepackage{fontspec}
\setmainfont{Latin Modern Roman}
\usepackage{amsmath,amssymb,mathtools}
\usepackage{xcolor,enumitem,booktabs,longtable,array,xurl,hyperref}
\definecolor{muted}{HTML}{53636C}
\hypersetup{colorlinks=true,urlcolor=blue,linkcolor=blue}
\setlength{\parindent}{0pt}
\setlength{\parskip}{5pt}
\newcommand{\R}{\mathbb R}
\newcommand{\E}{\mathbb E}
\newcommand{\N}{\mathbb N}
\newcommand{\ind}{\mathbf 1}
\newcommand{\Var}{\operatorname{Var}}
\newcommand{\Cov}{\operatorname{Cov}}
\newcommand{\supp}{\operatorname{supp}}
\DeclareMathOperator*{\argmin}{arg\,min}
\DeclareMathOperator*{\argmax}{arg\,max}
\newcommand{\entrymeta}[3]{{\sffamily\small\color{muted}\textbf{\texttt{#1}}\quad|\quad #2\par #3\par}}
\newcommand{\Problemref}[1]{\texttt{#1}}
\newcommand{\JELref}[2]{\texttt{#2} (JEL \texttt{#1})}
\begin{document}
\section*{Global financial safety net}
\textbf{Problem ID:} \texttt{F33-118}\quad \textbf{JEL:} \texttt{F33}\par
% BEGIN ECONOMICS STATEMENT
\label{problem:OP-0339}
{\sffamily\small\color{muted}Original subject group: International finance and exchange rates\par}
\label{definitions:problem:30007621}
\textbf{Audit classification:} Published research agenda. The official agenda supports global safety-net reform and tool combinations. The weighted architecture objective, access-cost and budget specification are editorial.
\subsection*{Definitions and precise research target}
Consider countries choosing reserve buffers and accessing central-bank swaps, regional credit facilities and multilateral lending. A safety-net architecture \(a\) specifies eligibility, committed capacities, prices, collateral, conditionality and activation rules. Countries choose prior borrowing and self-insurance anticipating access; lenders face a specified resource constraint. Let \(\omega\) be a shock state with a specified probability law. Let \(L_c(a,\omega)\) be real crisis losses in country \(c\), \(C_c(a)\) reserve and facility costs, and \(F(a,\omega)\) shared fiscal resource losses. Given nonnegative welfare weights \(w_c\), seek a feasible architecture minimizing
\[\sum_cw_c\{E[L_c(a,\omega)]+C_c(a)\}+E[F(a,\omega)].\]
The task is to characterize coverage, coordination and instrument substitution that improve insurance without inducing excessive prior risk. Define stigma as a measurable private access cost or signaling consequence, and incorporate it in country behavior. Account for correlated global demands for liquidity and currency-specific funding needs. Identify causal effects where past support was selected in response to severity. Report sensitivity to unequal eligibility and alternative welfare weights. Success in protecting eligible countries is insufficient to establish optimal global coverage or resource allocation.

\textbf{Additional formal definitions and scope (editorial).} Specify a finite country set, committed foreign-currency budgets, and admissible activation histories. Reserve costs include opportunity or resource costs with a fixed safe-rate benchmark; avoid counting asset purchases themselves as lost resources. Facility terms distinguish committed lending capacity from actual contingent disbursement and expected recovery. Country choices include borrowing, reserve holding and take-up under anticipated access. Stigma is a utility/resource cost or an inference affecting subsequent funding, with a stated signal model if used. The weighted objective fixes \(w_c\geq0\), usually normalized to sum to one, and puts every cost in common consumption units. Shared fiscal losses must not duplicate country-level losses. A feasible architecture is optimal conditional on these weights and commitments; global coverage under commitment and sequentially credible discretionary assistance are distinct policy problems.
\subsection*{Variants and assumption changes}
\begin{enumerate}
\item \textbf{Committed liquidity insurance (editorial).} Fix eligibility and budgets before countries choose reserves/borrowing; optimize common-currency crisis losses plus resource costs.
\item \textbf{Correlated shocks and selective access (editorial).} Allow global funding stress, stigma and discretionary activation; compare coverage and credible allocation under alternative country weights.
\end{enumerate}
\subsection*{References and source audit}
\begin{enumerate}
\item Official January 9, 2026 web version; locator: Interconnected World / International finance, safety-net paragraph. Broad agenda; exact model editorial. URL: \url{https://www.bankofengland.co.uk/research/bank-of-england-agenda-for-research}.
\end{enumerate}
\begin{enumerate}\raggedright
\item\hypertarget{definitions:source:30007621:1}{} Bank of England. 2026. \emph{Bank of England Agenda for Research, 2025–2028}. Interconnected World / International finance, safety-net paragraph.
\url{https://www.bankofengland.co.uk/research/bank-of-england-agenda-for-research}.
\end{enumerate}

\subsection*{Common conventions: Source mathematical conventions}

These conventions apply to each entry unless explicitly replaced.

\textbf{Models and identification.} A primitive model \(m\) specifies all probability laws, feasible choices, information, equilibrium and measurement rules used by its target. The admissible class \(\mathcal M\) is fixed independently of outcomes. If \(P_O(m)\) is its observable probability law and \(\tau(m)\) the target, its identified set at observable law \(P\) is
\[
 \mathcal I_\tau(P)=\{\tau(m):m\in\mathcal M,\ P_O(m)=P\}.
\]
Point identification means that this set is a singleton. An empty set signals incompatibility of data and assumptions. Coordinate bounds need not describe the joint set. Bounds are sharp only if every retained target is attainable or arbitrarily approachable by compatible models. Finite bounds require explicit restrictions on unobserved quantities; unrestricted confounding, hidden wealth, extrapolation or arbitrary equilibrium selection can yield uninformative sets.

\textbf{Causal interventions.} A potential outcome \(Y_i(p)\) is the outcome under a complete policy regime \(p\), including timing, financing, information and market spillovers. The target population and units are fixed before treatment. With interference, use the complete assignment vector or a justified exposure mapping; individual no-interference notation alone cannot identify an economy-wide intervention. A difference of mean potential outcomes does not require the cross-world joint distribution. Conditioning on post-treatment migration, survival, enrollment or employment changes the estimand and needs separate justification.

\textbf{Statistical statements.} An estimator, confidence set, forecast or test specifies a sampling law, model class and loss/coverage criterion. Uniform \((1-\alpha)\)-coverage means \(\inf_{m\in\mathcal M}\Pr_m\{\tau(m)\in C_n\}\geq1-\alpha\), or an explicitly asymptotic version. The whole parameter space trivially has coverage, so informativeness requires a stated width, risk or power objective. A forecasting improvement is relative to a named benchmark, information set and evaluation population.

\textbf{Equilibrium and optimization.} An equilibrium comprises feasible plans, optimality, beliefs where needed, budget constraints and all market/resource clearing. The primitives or a stated benchmark define each correspondence \(\mathcal E(p;m)\). Nonemptiness is assumed or proved, never inferred just from compact policy sets. A selected-equilibrium target uses a declared selection rule; a robust target quantifies over all permitted equilibria. Use suprema or infima unless attainment is established. An \(\varepsilon\)-optimal policy is feasible and within \(\varepsilon\) of the finite optimum value. Report an empty feasible set separately.

\textbf{Welfare and accounting.} Normative utility, interpersonal normalization, social weights, numeraire, discounting and the use of public receipts are primitives. Behavioral utility need not coincide with the welfare criterion. Household welfare includes ownership income and rents once; adding profits again to compensating variation generally double-counts them. A monetary equivalent is defined by an explicit transfer and price convention, with monotonicity and range conditions. Average effects, mechanism contributions, transfers and welfare losses are different targets. Infinite sums require convergence and dynamic feasible plans require terminal or no-Ponzi conditions.

\textbf{Source versions.} A working-paper date, revision date and journal publication year are distinguished. A DOI for a journal version is not automatically the DOI of an earlier manuscript. Locators refer to the stated version and printed pages where specified. A metadata-only check does not verify a claim about a full-text conclusion. Every entry's source subsection narrows the provenance claim accordingly.
% END ECONOMICS STATEMENT
\end{document}
